\chapter{Experimental Design and Training Objectives} \label{ch:method}
This chapter describes the formal problem definitions, dataset characteristics, and the specific configurations of the research objectives.

% ================================================================
\section{Problem Formulation} \label{sec:method:problem}
% ================================================================
\begin{planned}
    \item Mathematical notation of the continuous glucose input sequence and the multi-step output trajectory vector.
    \item Formal presentation of the binary event classification label representing hypoglycemia breaches.
\end{planned}

% ================================================================
\section{Dataset and CGM Characteristics} \label{sec:method:setup}
% ================================================================
\subsection{The T1DATA Cohort} \label{sec:method:setup:cgm}
\begin{planned}
    \item Introduction of the T1DATA clinical trial dataset: 36 anonymised Type~1 diabetes patients, CGM readings at a 5-minute sampling interval.
    \item The glucose signal has a lag-1 autocorrelation above 0.98, meaning each reading is nearly identical to the previous one. This makes it easy for a model to predict well by copying recent history, which is the root cause of the time-shift artefact.
    \item Class imbalance: only 5.7\% of all labelled windows contain a hypoglycaemia event, motivating the F2 metric and weighted loss in direct classifiers (see also Fig.~\ref{fig:bg:class_imbalance}).
\end{planned}

\subsection{Train/Val/Test Splits and Hyperparameter Optimization} \label{sec:method:setup:splits}
\begin{planned}
    \item Chronological partition (60/20/20) to preserve continuous temporal integrity and prevent future-data leakage.
    \item Operational parameters of the Optuna automated search space framework executed exclusively on the designated tuning patient.
\end{planned}


% ================================================================
\section{Training Objectives} \label{sec:method:objectives}
% ================================================================
\begin{planned}
    \item \textbf{Objective 1 (Offset-Aware Loss):} Both models are retrained twice --- once with the Bounded-Lag loss and once with Soft-DTW --- while all other settings remain unchanged. The goal is to test whether using a loss that tolerates small timing differences can reduce the prediction delay.
    \item \textbf{Objective 2 (Multi-Step):} The single 60-minute output is replaced by 12 simultaneous step predictions covering the full forecast window. Supervising all steps at once gives the model direct feedback about earlier time points, which may anchor the prediction in time.
    \item \textbf{Objective 3 (Event-Centric):} Two approaches are tested: (a) the tolerance window $\tau$ is swept from 5 to 60 minutes to see how much loosening the timing requirement improves detection for the MSE-trained forecasters; (b) direct binary classifiers are trained end-to-end on the event label and compared against the threshold-based forecast detectors.
\end{planned}

\begin{table}[htbp]
\centering
\caption{Hyperparameters derived by Optuna (50 trials on patient 85202) and locked for all experiments. These values are held constant across all three research objectives; only the loss function or output head is varied per objective.}
\label{tab:method:hyperparams}
\begin{tabular}{lrr}
\toprule
Hyperparameter & LSTM & PatchTST \\
\midrule
Hidden size / $d_\text{model}$  & 256   & 64 \\
Number of layers                & 2     & 3 \\
Dropout                         & 0.159 & 0.025 \\
Learning rate                   & $9.61 \times 10^{-4}$ & $1.21 \times 10^{-4}$ \\
Batch size                      & 128   & 128 \\
Feed-forward dim.\ ($d_\text{ff}$) & ---   & 256 \\
Attention heads ($h$)           & ---   & 4 \\
Lookback window                 & \multicolumn{2}{c}{24 steps (120 min)} \\
Forecast horizon                & \multicolumn{2}{c}{12 steps (60 min)} \\
Optuna trials                   & \multicolumn{2}{c}{50} \\
\bottomrule
\end{tabular}
\end{table}
